% Общие методы обработки данных минископного кальциевого имиджинга.
% Описаны один раз; на них ссылаются секции 2.1 (гиппокамп), 2.4 (размерности), 3.1 (INTENSE).

\section{Методы обработки данных минископного кальциевого имиджинга}\label{sec:calcium_methods}

Данный раздел описывает единый конвейер обработки записей однофотонного минископного кальциевого имиджинга, применявшийся во всех экспериментах настоящей работы.

\subsection{Экстракция нейронных компонент}\label{sec:cnmfe_processing}

Данные кальциевого имиджинга обрабатывались с помощью BEARMIND --- специализированного конвейера анализа, разработанного авторами (см.\ раздел~\ref{sec:driada}), построенного на основе фреймворка CaImAn~\cite{Pnevmatikakis2016}. Программное обеспечение предназначено для сквозной обработки многосессионных записей минископа.

Каждая сессия начиналась с визуального осмотра для определения границ обрезки. Предобработка включала коррекцию фона и удаление артефактов движения алгоритмом NoRMCorre~\cite{Pnevmatikakis2017}. Затем видеозаписи разбивались на пространственные маски и временные ряды активности методом CNMF-E~\cite{Zhou2018}. Гиперпараметры CNMF-E подбирались с использованием заданных критериев качества (компактность пространственной маски, SNR временного ряда и структура остатков) и верифицировались визуальным осмотром. Основные диапазоны параметров:
\begin{itemize}
    \item \texttt{gSig} (оценка размера нейрона): 3--6 пикселей (при разрешении 3~мкм/пиксель);
    \item \texttt{min-corr} (порог локальной корреляции): 0{,}85--0{,}95;
    \item \texttt{min-pnr} (отношение пикового сигнала к шуму): 3--15.
\end{itemize}

Ключевые результаты были качественно устойчивы в данном диапазоне параметров (анализ робастности).

После декомпозиции выделенные компоненты предварительно фильтровались по порогам \texttt{rval-thr} = 0{,}95 и \texttt{min-SNR} = 3, в соответствии с рекомендациями CaImAn. Финальная классификация компонент на нейроны и артефакты выполнялась с помощью автоматизированной системы контроля качества, описанной в следующем разделе.

Сигналы флуоресценции нормализовались как $dF/F$. Пространственные карты регистрировались между сессиями с помощью CellReg~\cite{Sheintuch2017} для продольного отслеживания индивидуальных нейронов.

% Драфт: Автоматизированный контроль качества компонент кальциевого имиджинга
% Источник: autoinspect_report_ru.docx
% Куда: после «Обработка данных кальциевого имиджинга» в update_nof.tex

\subsection{Автоматизированный контроль качества выделенных компонент}\label{sec:autoinspect}

Как отмечалось выше, после декомпозиции видеозаписей алгоритмом CNMF-E полученные компоненты необходимо классифицировать на нейроны и артефакты. Стандартная процедура предполагает ручную верификацию каждого кандидата экспертом, что требует значительных временных затрат и приводит к межэкспертной вариабельности оценок. Для автоматизации этого этапа была разработана система машинного обучения, обучающаяся на решениях эксперта и воспроизводящая их с высокой точностью.

\paragraph{Пространство признаков.}
Для каждого кандидата в нейроны вычисляется 35 признаков, охватывающих пять категорий:
\begin{itemize}
    \itemsep0em
    \item \textit{Пространственные и морфологические признаки} (11): площадь, выпуклость, соотношение сторон, компактность, эксцентриситет и другие характеристики пространственного следа.
    \item \textit{Временные и статистические признаки} (9): базовая флуоресценция, дисперсия, коэффициенты асимметрии и эксцесса сигнала $dF/F$, автокорреляционные характеристики.
    \item \textit{Событийные признаки} (8): частота кальциевых событий, событийное отношение сигнал/шум, качество аппроксимации при моделировании кинетики кальция, времена нарастания и спада.
    \item \textit{Метрики качества реконструкции} (4): оценка соответствия наблюдаемого сигнала генеративной модели кальциевой динамики.
    \item \textit{Алгоритмические метрики CaImAn} (2): встроенные оценки SNR и пространственной корреляции.
\end{itemize}

Детекция кальциевых событий, необходимая для вычисления событийных признаков, выполняется вейвлет-методом из пакета DRIADA (см.\ раздел~\ref{sec:driada}). Для компонент со сложно характеризуемыми паттернами активности реализована иерархическая система из пяти уровней детекции с последовательным ослаблением критериев качества, обеспечивающая покрытие свыше 99\% кандидатов.

\paragraph{Классификатор.}
В качестве модели классификации использовались объяснимые бустинг-машины (Explainable Boosting Machines, EBM)~\cite{InterpretML-EBM}~--- обобщённые аддитивные модели с парными взаимодействиями~\cite{Lou2013-GA2M}, в которых вклад каждого признака представлен индивидуальной функцией формы:
\begin{equation}\label{eq:ebm}
    g(y) = \beta_0 + \sum_{j=1}^{p} f_j(x_j) + \sum_{(i,j)} f_{ij}(x_i, x_j),
\end{equation}
где $f_j$~--- одномерные функции формы, а $f_{ij}$~--- парные взаимодействия. Такая структура обеспечивает интерпретируемость модели: для каждой классификации можно определить вклад каждого признака в итоговое решение. При этом производительность EBM сопоставима с ансамблевыми методами (случайные леса, градиентный бустинг).

Анализ важности признаков показал, что наибольшую дискриминационную способность обеспечивают временная статистика сигналов (дисперсия, эксцесс) и событийные характеристики (событийный SNR), тогда как встроенные метрики CaImAn, традиционно используемые для фильтрации, оказались менее информативными.

\paragraph{Валидация.}
Система оценивалась методом перекрёстной валидации на наборе из 92\,242 размеченных компонент (187 сессий записи, 4 эксперимента). С учётом асимметрии ошибок классификации (ложноположительные результаты~--- сохранение артефактов~--- более критичны, чем пропуск нейронов) порог принятия решения оптимизировался по метрике $F_\beta$ с $\beta = \sqrt{1/3}$, придающей повышенный вес точности. Достигнутая производительность: точность (precision) 97{,}6\%, полнота (recall) 93{,}3\%, общая метрика $F_\beta = 96{,}5\%$.

Автоматизированная система интегрирована в конвейер BEARMIND и позволяет сократить время контроля качества с нескольких часов до нескольких минут на сессию, обеспечивая при этом воспроизводимость решений.


\subsection{Детекция кальциевых событий}\label{sec:event_detection}

Для извлечения дискретных кальциевых транзиентов из сигналов $dF/F$ использовался вейвлет-метод, ранее описанный в~\cite{Neugornet2021}, с незначительными модификациями. Данный метод был выбран за его устойчивость к шуму и способность обнаруживать мультимасштабные особенности в данных флуоресценции.

Сигнал $dF/F$ предварительно сглаживался свёрткой с гауссовым ядром ($\sigma = 8$), затем преобразовывался с помощью непрерывного вейвлет-преобразования (CWT), что давало двумерный массив вейвлет-коэффициентов по времени и масштабу. Значимые события проявлялись как <<гребни>> локальных максимумов на соседних масштабах.

Использовались обобщённые вейвлеты Морса~\cite{Lilly2012} $\Psi_{\beta,\gamma}(\omega) = \alpha_{\beta,\gamma}\omega^{\beta}e^{-\omega^{\gamma}}$ с параметрами $\gamma = 3$ и $\beta = 2$, соответствующими так называемым вейвлетам Эйри, которые обеспечивают минимальную площадь Гейзенберга и хорошую временнýю локализацию.

Для каждого масштаба определялись локальные максимумы вейвлет-амплитуды. Максимумы на наибольшем масштабе инициировали гребни, которые затем продлевались на меньшие масштабы по критерию временнóго совпадения. При наличии нескольких кандидатов на данном шаге сохранялся только наивысший локальный максимум в пределах каждого гребня, остальные инициировали новые гребни. Процесс продолжался до наименьшего масштаба, после чего гребни, не охватывающие все масштабы, отбрасывались.

После построения всех гребней применялся фильтр значимости по длительности гребня, пиковой амплитуде и масштабу, на котором достигался максимум. Гребни, прошедшие фильтрацию, классифицировались как кальциевые события.
